\chapteropen{02}

\begin{cele}
  \item podać, w jakim przedziale mieści się typowe zapotrzebowanie dorosłych na sen i dlaczego jest to przedział, a nie cel;
  \item odróżnić potrzebę snu, nawyk snu i dług snu oraz przewidzieć, co każde z nich robi z planem na 12 godzin;
  \item rozpoznać sytuacje, w których długie spanie jest objawem wymagającym konsultacji lekarskiej, i wiedzieć, o co zapytać;
  \item zmierzyć własne zapotrzebowanie na tyle uczciwie, by wynik nie był tylko wrażeniem;
  \item rozróżnić czas spędzony w łóżku od czasu snu i policzyć wydajność snu na własnych danych.
\end{cele}

\begin{uwaga}[Ta książka nie jest poradą medyczną]
  Ten rozdział nie diagnozuje, nie leczy i nie zastępuje lekarza, psychologa ani specjalisty snu. Podaje ogólne zakresy, sposoby obserwacji i listę sytuacji, w których rozmowa z lekarzem jest zasadna. Jeśli coś z tej listy dotyczy Ciebie, nie szukaj odpowiedzi w książce, tylko umów się na wizytę. Autor nie jest lekarzem. Nie znasz swojej historii zdrowia, a książka nie ma jej przed sobą.
\end{uwaga}

\section{Ile snu zwykle potrzebuje dorosły}

Zanim policzysz cokolwiek w planie dnia, potrzebujesz punktu odniesienia. Najczęściej przywoływany zakres dla dorosłych to \verify{około 7--9 godzin snu na dobę}, a niżej ludzie, u których wyraźnie gorzej działa pamięć, uwaga i nastrój, znajdują się zwykle w grupie, która śpi krócej niż 7 godzin. Ten zakres jest opisem populacji, a nie celem, który każdy ma osiągnąć co noc.

To ważne dla planu z 12 godzinami, bo 12 godzin leży wyraźnie poza tym zakresem. Nie oznacza to automatycznie, że coś jest nie tak. Oznacza tylko, że osoba, która chce spać 12 godzin, musi umieć wyjaśnić, dlaczego właśnie ona potrzebuje więcej niż większość, i sprawdzić, czy to wyjaśnienie wytrzymuje obserwację. Wielu ludzi zakłada, że ich potrzeba jest wyjątkowa, zanim sprawdzi cokolwiek, co ją opisuje.

Dwie rzeczy trzeba rozdzielić na samym początku. Pierwsza to \emph{ile snu organizm potrzebuje, żeby funkcjonować w dzień}. Druga to \emph{ile snu chcesz mieć}. Obie wartości mogą się różnić, a w praktyce różnią się często. Książka dotyczy obu, ale rozdział ten zajmuje się przede wszystkim pierwszą, bo tylko ona mówi cokolwiek o zdrowiu i o tym, czy plan jest bezpieczny.

\section{Różnice indywidualne i to, czego nie wiemy}

Ludzie różnią się zapotrzebowaniem na sen, i to nie tylko w sensie upodobań. Część tej różnicy wiąże się z genetyką i z tym, kiedy organizm chce zasypiać i wstawać, czyli z chronotypem. Część wiąże się z wiekiem, z poziomem aktywności fizycznej, z tym, co się je i pije, oraz z tym, jakie leki się przyjmuje. \verify{Uwarunkowania genetyczne długości snu są przedmiotem badań, ale dokładny udział każdego czynnika w zmienności zapotrzebowania jest nadal przedmiotem dyskusji}.

Praktyczny wniosek jest prosty i trochę niewygodny: nie ma jednego testu, który w kilka minut powie Ci, ile śpisz optymalnie. Samo poczucie, że \glqq wyspałem się\grqq{}, bywa mylące. Osoba, która od lat śpi krótko i przywykła do stałego lekkiego zmęczenia, może uznać je za normę. Osoba, która zawsze śpi długo, może uznać za normę to, co jest tylko skutkiem nieregularnego dnia.

Dlatego ten rozdział nie podaje Ci liczby do przyjęcia. Podaje sposób, w jaki możesz zebrać dane o sobie i ocenić je bez wiary w intuicję. Sekcja o pomiarze jest najważniejszą częścią rozdziału, nawet jeśli wygląda najmniej efektownie.

\section{Dług snu i jego rachunek}

Dług snu to różnica między tym, ile spałeś, a tym, ile organizm faktycznie potrzebował. Działa jak każdy dług: rośnie, gdy wydajesz więcej, niż masz. Rachunek jest prosty, bo ma tylko jedno działanie.

Jeśli Twoje zapotrzebowanie to 8 godzin, a śpisz 7 godzin, to każdej nocy brakuje godziny. Przez siedem nocy daje to 7 godzin w tygodniu. Jeśli śpisz 6 godzin, deficyt to 2 godziny na noc i 14 godzin w tygodniu. Ta arytmetyka jest pewna. Nie jest pewne, jak ten deficyt przekłada się na funkcjonowanie u konkretnej osoby, i nie jest pewne, czy weekend go skasuje.

\verify{Badania nad skutkami ograniczonego snu pokazują pogorszenie uwagi i reakcji po kilku dniach krótszego snu, a odzyskiwanie tej sprawności w weekend bywa tylko częściowe}. Dla tej książki wystarcza prosty wniosek: dług snu nie znika sam, a sen w weekend nie jest zerowaniem licznika, tylko częściową spłatą.

Ma to znaczenie dla planu z 12 godzinami w obie strony. Ktoś, kto ma w tygodniu dni krótkiego snu i dni długiego, nie ma dwóch różnych rytmów, tylko jeden rytm z nieregularnością. Nieregularność sama w sobie utrudnia ocenę, czego naprawdę potrzebujesz. Dlatego pomiar, który zobaczysz za chwilę, wymaga stałej pory zasypiania przez co najmniej kilka dni przed tym, jak zaczniesz wyciągać wnioski.

\section{Kiedy długi sen to powód do rozmowy z lekarzem}

Są sytuacje, w których dłuższy sen nie jest wyborem stylu życia, tylko jednym z objawów czegoś, co warto zbadać. Poniższa lista nie jest diagnozą. Jest listą pytań, które warto zadać lekarzowi, jeśli dotyczą Ciebie.

\begin{itemize}
  \item \textbf{Nadmierna senność w ciągu dnia}, czyli zasypianie w sytuacjach, w których nie powinieneś, na przykład w trakcie rozmowy, za kierownicą albo w ciągu dnia mimo odpowiedniej liczby godzin w nocy. \verify{Senność dzienna, która utrudnia codzienne funkcjonowanie, jest jednym z głównych powodów konsultacji w zaburzeniach snu}.
  \item \textbf{Chrapanie z przerwami w oddychaniu}, zauważone przez partnera albo domownika. \verify{Bezdech senny jest zaburzeniem, które często zostaje niezauważone przez osobę, której dotyczy, a sygnalizuje je przede wszystkim otoczenie}.
  \item \textbf{Obniżony nastrój, który trwa tygodniami}, zwłaszcza gdy towarzyszy mu utrata zainteresowań albo zmęczenie, które nie mija po śnie. \verify{Zaburzenia nastroju mogą zmieniać ilość snu i jakość snu w obu kierunkach, a sam sen nie zawsze wystarcza do ich rozróżnienia}.
  \item \textbf{Objawy sugerujące zaburzenia tarczycy} lub inne zaburzenia metaboliczne, na przykład wyraźny spadek energii, zmianę masy ciała bez zmiany diety albo zmianę wrażliwości na zimno. \verify{Niedoczynność tarczycy może obniżać energię i zwiększać zapotrzebowanie na sen, ale te same objawy mają wiele innych przyczyn}.
  \item \textbf{Nowe leki albo zmiana dawki}, po których sen wyraźnie się zmienił. Wiele leków wpływa na sen, a lekarz lub farmaceuta może powiedzieć, czy zmiana jest oczekiwana.
\end{itemize}

Warto zauważyć, czego na tej liście nie ma: samego faktu, że śpisz długo. Jeśli śpisz 12 godzin, czujesz się wypoczęty, nie masz senności w ciągu dnia, nie chrapiesz z przerwami w oddychaniu i nie masz innych objawów, to sam czas snu nie jest powodem do alarmu. Jest natomiast powodem do uczciwej obserwacji i do zapisywania danych, bo długi sen, który zaczyna się z czasem wydłużać bez wyraźnej przyczyny, jest sygnałem wartym sprawdzenia.

\begin{uwaga}[Czego nie robić z tą listą]
  Nie wybieraj się do lekarza z gotową diagnozą wyczytaną z internetu i nie zmieniaj leków na własną rękę, bo ktoś w książce napisał, że \glqq to tarczyca\grqq{}. Przynieś dziennik snu z sekcji o pomiarze, opisz objawy konkretnie i pozwól lekarzowi zadać pytania, których nie przewidziałeś.
\end{uwaga}

\section{Jak zmierzyć własne zapotrzebowanie}

Pomiar wymaga trzech decyzji: ile czasu, co zapisywać i jak nie zepsuć wyniku. Najczęstsza pomyłka polega na tym, że ludzie robią pomiar w czasie, gdy wszystko jest wyjątkowe: w wakacje, w czasie choroby, tuż przed ważnym terminem albo po zarwanej nocy. Wynik z takiego okresu mówi o tym okresie, a nie o Tobie.

Zalecam \emph{tydzień bez budzika} jako minimum, a dwa tygodnie jako wygodniejszą wersję. Przez te dni nie ustawiasz budzika i nie wstajesz na czas, o ile nie musisz. Jeśli musisz, zapisz to. Pierwsze dni po wcześniejszym niedospaniu mogą być zakłócone, dlatego pierwszy tydzień często nie jest najbardziej wiarygodny. Ten wybór jest rozsądnym kompromisem między czasem a wiarygodnością, a nie normą ustaloną przez badania.

Co zapisywać każdego dnia:

\begin{enumerate}
  \item godzinę, o której położyłeś się do łóżka;
  \item godzinę, o której zgasiłeś światło;
  \item szacunkowy czas zaśnięcia, czyli ile minut minęło od zgaszenia światła do momentu, w którym przestałeś być świadomy;
  \item liczbę przebudzeń w nocy i ich długość, jeśli je pamiętasz;
  \item godzinę, o której wstałeś, oraz godzinę, o której faktycznie wyłączyłeś budzik albo pierwszy raz wstałeś z łóżka;
  \item ocenę samopoczucia w ciągu dnia w skali od 1 do 5, zapisaną raz, np. około południa.
\end{enumerate}

Punkty od pierwszego do piątego pozwalają policzyć dwie rzeczy, które warto odróżniać. Pierwsza to czas w łóżku, czyli od położenia się do wstania. Druga to czas snu, czyli czas, w którym faktycznie spałeś. Różnica między nimi to czas leżenia bez snu, czyli czas zasypiania, przebudzenia i przebywania w łóżku bez snu. Ten podział jest najważniejszym narzędziem tego rozdziału.

\begin{przyklad}[Dwie liczby zamiast jednej]
  Osoba prowadzi dziennik przez tydzień. Średnio kładzie się o 22:30 i wstaje o 10:30, więc czas w łóżku to 12 godzin. Zapisuje, że zasypia po około 40 minutach, budzi się trzy razy po kilka minut i wstaje z łóżka około 10:30. Z zapisów wynika, że czas snu to około 10 godzin 20 minut, a nie 12. W tym przykładzie 12 godzin w łóżku to tylko 10 godzin snu. Dopiero ta różnica pokazuje, co faktycznie dzieje się w nocy, i jest punktem wyjścia do dalszej rozmowy, na przykład z lekarzem.
\end{przyklad}

\section{Dwanaście godzin w łóżku a dwanaście godzin snu}

Najczęstszy błąd w planie z długim snem to założenie, że czas w łóżku jest czasem snu. Jest to wygodne, bo prowadzi do prostej arytmetyki: 12 godzin w łóżku równa się 12 godzinom odpoczynku. Arytmetyka jest poprawna, ale założenie nie zawsze jest prawdziwe.

Wydajność snu to stosunek czasu snu do czasu w łóżku wyrażony w procentach. Jeśli spędzasz w łóżku 12 godzin, a śpisz 8, wydajność wynosi $8/12 \approx 67\%$. Jeśli spędzasz 12 godzin i śpisz 10 godzin 30 minut, wydajność wynosi $10{,}5/12 = 87{,}5\%$. Obie liczby opisują tę samą liczbę godzin w łóżku, a mówią o dwóch różnych nocach.

\begin{falle}[Skoro śpię 12 godzin, to odpoczywam 12 godzin]
  Nie wiadomo, ile z tych godzin było snem. Leżenie w łóżku z zamkniętymi oczami nie jest równoważne ze snem, a czas w łóżku obejmuje też zasypianie, przebudzenia i czas, gdy się przewracasz i myślisz. Jeśli plan opiera się na 12 godzinach w łóżku, a nie na 12 godzinach snu, rachunek jest zawyżony. Zanim zbudujesz wokół tego harmonogram, policz wydajność z własnych danych.
\end{falle}

Dla planu dnia ma to konkretne konsekwencje. Jeśli śpisz 12 godzin w łóżku, ale tylko 8 godzin rzeczywiście, to przez siedem nocy w tygodniu odzyskujesz tylko 56 godzin snu, a zajmujesz łóżko przez 84 godziny. Różnica 28 godzin to czas, który w planie wygląda jak sen, a w rzeczywistości jest leżeniem. Może to być w porządku, jeśli leżenie jest zamierzone, ale nie powinno się udawać, że jest snem.

\section{Sen a wiek}

Zapotrzebowanie na sen zmienia się w ciągu życia, a zmiana nie jest liniowa. \verify{U niemowląt i małych dzieci łączna długość snu jest znacznie większa niż u dorosłych, u nastolatków przesuwa się pora zasypiania, a u osób starszych częściej pojawia się sen krótszy i bardziej przerywany}. Te ogólne tendencje są dobrze opisane, ale pojedyncza osoba może od nich odbiegać, i to bez powodu do niepokoju.

Dla czytelnika, który ma 35 do 50 lat, wniosek jest taki, że wiek sam w sobie nie tłumaczy długiego snu. Jeśli 12 godzin w łóżku pojawiło się nagle, po latach krótszego snu, to jest to zmiana, którą warto zbadać. Jeśli 12 godzin jest stałym rytmem od młodości i nie towarzyszą mu objawy z wcześniejszej listy, to jest to zwykle sprawa indywidualnej potrzeby, którą trzeba jednak zmierzyć, a nie założyć.

Osobną sprawą jest to, jak wiek wpływa na samą ocenę. Im dłużej ktoś żyje z pewnym rytmem, tym mniej zauważa, że ten rytm jest zmęczony. Dlatego pomiar wykonany w spokojnym okresie jest bardziej wiarygodny niż ocena z pamięci.

\section{Jak czytać swoje dane}

Po dwóch tygodniach zapisów masz czternaście wierszy. Policz z nich trzy wartości: medianę czasu snu, medianę wydajności i liczbę dni, w których samopoczucie było wyraźnie gorsze. Mediana jest lepsza od średniej, bo jedna bardzo długa noc nie przesuwa wyniku do góry.

Następnie porównaj samopoczucie z czasem snu. Jeśli dni z dłuższym snem nie są lepsze, długi sen nie jest lekarstwem na zmęczenie, tylko jego towarzyszem. Jeśli dni z dłuższym snem są wyraźnie lepsze, a Ty nie masz objawów z listy, to masz dane, które możesz pokazać lekarzowi lub wykorzystać w planie. W obu przypadkach wynik jest tylko tym, co zapisałeś, a nie tym, co chciałbyś zobaczyć.

Jedno zastrzeżenie. Dziennik snu nie jest badaniem. Nie mierzy fal mózgowych ani tego, jak głęboki był sen. Jest narzędziem do obserwacji własnego rytmu, a nie do stawiania diagnoz. Jego wartość polega na tym, że jest tani, uczciwy i zrobiony przez Ciebie, a nie przez wrażenie zapamiętane po tygodniu.

\begin{cwiczenia}
  \exitem{Oblicz, ile godzin tygodniowo różni się czas w łóżku od 12 godzin dziennie a czas w łóżku 8 godzin dziennie, przy siedmiu nocach.}{Różnica dzienna to 4 godziny (12 minus 8). Przy siedmiu nocach daje to 28 godzin tygodniowo ($4 \cdot 7 = 28$).}
  \exitem{Osoba leży w łóżku przez 12 godzin i zapisuje czas snu 7 godzin 20 minut. Jaka jest wydajność snu?}{Czas snu to 440 minut, czas w łóżku to 720 minut. Wydajność wynosi $440/720 \approx 61\%$.}
  \exitem{Ktoś śpi 6 godzin przez pięć kolejnych nocy, a w weekend 8 godzin dziennie. Ile godzin snu brakuje w tygodniu względem 8 godzin każdej nocy, jeśli weekend nie zmienia niedoboru?}{Pięć nocy po 6 godzin daje 30 godzin, a dwie noce po 8 godzin dają 16, więc łącznie 46 godzin. Pełne 56 godzin minus 46 daje 10 godzin brakujących. Ta odpowiedź zakłada, że weekend tylko częściowo spłaca dług, bo wynik jest sumą faktycznego snu, nie spłatą.}
  \exitem{Wymień trzy objawy, które uzasadniają rozmowę z lekarzem, nawet jeśli czas snu wydaje się wystarczający.}{Na przykład: nadmierna senność w ciągu dnia (zasypianie w sytuacjach, w których nie powinno się to zdarzać), chrapanie z przerwami w oddychaniu zgłaszane przez partnera albo domownika, obniżony nastrój utrzymujący się przez tygodnie.}
  \exitem{Jakie sześć rzeczy warto zapisać codziennie, żeby odróżnić czas w łóżku od czasu snu?}{Godzinę położenia się do łóżka, godzinę zgaszenia światła, szacowany czas zaśnięcia, liczbę i długość przebudzeń, godzinę wstania i samopoczucie w skali od 1 do 5 zapisane raz w ciągu dnia.}
  \exitem{Dlaczego pomiar w czasie wakacji albo po chorobie nie jest dobrym punktem odniesienia dla własnego zapotrzebowania?}{Bo w takim okresie zmienia się wiele naraz: rytm dnia, obciążenie, dług snu z wcześniejszych tygodni i stan zdrowia. Wynik opisuje ten okres, a nie typowy rytm osoby.}
  \exitem{Dlaczego mediana jest lepszym wskaźnikiem niż średnia w dzienniku snu?}{Bo pojedyncza bardzo długa noc, na przykład po zarwanej nocy, przesuwa średnią w górę i daje obraz rytmu, którego nie ma. Mediana ignoruje pojedyncze skrajne wartości i lepiej opisuje typową noc.}
\end{cwiczenia}

\zadanie{Przez 14 kolejnych dni, bez budzika tam, gdzie to możliwe, prowadź dziennik z sześcioma polami opisanymi w sekcji o pomiarze. Na koniec policz medianę czasu w łóżku, medianę czasu snu, medianę wydajności i liczbę dni z samopoczuciem 1 lub 2. Miara sukcesu: masz 14 kompletnych wierszy, wyliczone trzy mediany oraz jedno zdanie o tym, czy dni z dłuższym snem były lepsze. Jeśli w dzienniku pojawia się senność w ciągu dnia albo chrapanie z przerwami w oddychaniu, dopisz to do wyniku i pokaż zapisy lekarzowi.}
